\section{Comparison and the remaining long-time problem}
\label{sec:comparison}
\subsection{Stationary identities and uniform orbit segments}
The relation between stationary entry times and section return times is
classical. Marklof \cite{Marklof} develops this relation through stationary
point processes and Palm distributions, without an ergodicity assumption.
The suspension identity in Section~\ref{sec:flux} is an elementary version
adapted to a finite physical record. It supplies the preparation factor
and the residual-time integral, not the geometry or mass of its active set.
The new calculation starts with $S_j\tau_R$, rather than with a supplied
independent roof sequence.

The use of length as a billiard generating function, and the relation of
action Hessians to linearized dynamics, likewise belong to classical
variational mechanics. Bolotin and Treschev \cite{Hill} treat Hill-type
determinant identities for discrete and continuous Lagrangian systems,
including symplectic twist maps. Our corner-cofactor identity is a finite
Dirichlet calculation of this kind; it is not claimed as a new general
Hill formula. Its role here is to compare the exponentially small physical
flux with its quadratic value \emph{relatively}, uniformly in the number
of reflections. The trace-norm bound in \eqref{eq:log-det} and the summable
endpoint bridge make this possible on one fixed neighborhood.

For a fixed $j$, ordinary Morse integration would yield a quadratic onset
with a constant depending on that segment. It would not, by itself, give
a common collar for all $j$, relative remainder bounds after division by
$\sinh(j\chi_R)^{-1}$, or uniform bounds for additive marks and sampling
cuts. Lemmas~\ref{lem:localization}, \ref{lem:bridge} and
\ref{lem:relative} give these additional statements for the complete
mechanical event. Theorem~\ref{thm:marked} then gives the summable
source-dependent hierarchy, and Section~\ref{sec:records} identifies
its joint record. The comparison with an independent roof having the same
marginal shows a change of onset exponent, not merely a different constant.

\subsection{Instability from statistical thresholds and from marked lengths}
B\'alint, De Simoi, Kaloshin and Leguil \cite{BDKL} recover curvatures at
period-two points and Lyapunov exponents of periodic orbits from the marked
length spectrum of open dispersing billiards with three strictly convex
obstacles satisfying the non-eclipse condition. Their datum consists of
labelled periodic lengths. Our coefficient sequence comes instead from
full-equilibrium probabilities of maximal collision counts in a periodic
circular billiard. Formula~\eqref{eq:ratio} recovers the normal
orbit's multiplier from these statistical thresholds. This is a different
observable, not a claim to stronger inverse-geometric rigidity. The
identification of the multiplier with the Jacobi recurrence is proved
explicitly in Section~\ref{sec:consequences}.


\subsection{Periodic rare events and physical collision onsets}
Carney, Nicol and Zhang \cite[Theorems 1--2]{CNZ} prove extreme-value
and compound-Poisson limits for visits to shrinking neighborhoods of
regular periodic points, including finite- and infinite-horizon Sinai
billiard maps. Their metric is adapted to stable and unstable foliations;
the extremal index is $1-|DT^p|_{E^u}|^{-1}$, and the associated
cluster multiplicities are geometric. Their observation length diverges
while the target shrinks so that the expected exceedance count stays
finite. Thus statistical detection of periodic instability is already
present in that work; it is not an originality claim of this article.

Our observable is the maximal number of physical impacts in a window
near $jg_e$, under full flow preparation. The collar, relative error and
shape/source derivatives are uniform in $j$, rather than derived from
a rare-visit limit. Theorems~\ref{thm:g-stability} and
\ref{thm:g-marked} verify those claims under noncircular deformations
and at competing threshold surfaces. Neither a rare-visit law nor a
fixed-segment Morse asymptotic supplies this uniform physical statement.
The proof does not invoke a mixing rate. Conversely it does not prove
a compound-Poisson law for long observation sequences.

The inverse issue is likewise explicit. In the circular model the
threshold spacing already fixes the radius and multiplier. We retain
Corollary~\ref{cor:multiplier} as its exact amplitude identity, not as
recovery of another unknown degree of freedom. In
Theorem~\ref{thm:g-inverse}, three independent curvature values vary
while every threshold location remains fixed. The unlabelled amplitudes
resolve precisely this uncertainty, with the unknown preparation area
removed by recovery from the same data. This does not assert full shape
rigidity or a stability estimate for noisy inverse data. It gives a
positive geometric consequence not determined by the threshold locations.

The marked-length results in \cite{BDKL} concern labelled periodic-orbit
lengths in open dispersing billiards and recover dynamical and curvature
information from those data. Our equal-gap family and coefficient
inversion distinguish the particular full-equilibrium data used here;
they do not purport to supersede that inverse theory. The variational
relation of action Hessians and multipliers remains classical
\cite{Hill}. The new relative estimate is its uniformly controlled use
for a complete physical threshold event with deformable geometry.

\subsection{Perturbative limit laws and the original Fourier programme}
Demers and Zhang \cite{DemersZhang} construct a functional-analytic
framework for billiard-map perturbations, including scatterer motions and
deformations. Their perturbation estimates and continuity of spectral
data address singularities on a much larger dynamical phase space than
our onset charts. Theorem~\ref{thm:marked} does not replace those operator
estimates. It gives every fixed radius derivative of a specified extremal
record, including its threshold singularity, in a regime selected by an
exact physical event. It is not an assertion of all-order differentiability
of the full transfer operator under arbitrary deformations.

Demers, Melbourne and Nicol \cite{DMN} prove martingale approximation and
statistical limit laws for H\"older observables of the time-one map of a
finite-horizon planar periodic Lorentz gas. Those typical-orbit limit laws
and the present maximal-count asymptotics concern different regimes.
In particular, neither their stated central limit and invariance principles
nor our threshold series is a source-differentiated mixed local Edgeworth
expansion for general collision records.

The earlier arithmetic and Fourier-inversion chapter \cite{A2R33} proves
an algebraic separation estimate and integrates four explicitly assumed
frequency envelopes. A mechanical application requires the characteristic
family itself to satisfy those envelopes, together with a central
integrable remainder and the required lattice and covariance conditions.
The present results provide a uniform, source-dependent calculation on a
nontrivial family of actual multi-impact events. They do not estimate the
complementary itineraries or construct the stable-curve operators needed
for that original long-time application. In particular, the sum in
\eqref{eq:summable} ranges over distinct onset windows; it cannot be
substituted for the characteristic function of one unrestricted long
record. Establishing that substitution's actual prerequisites is a
separate mathematical problem, not a consequence of the notation.

Thus the structural conclusion is a collision-order-uniform geometric
threshold law, including noncircular stability, competing onsets,
curvature identification and specified marked response under full preparation. The
fixed-window theorem in \cite{TC}, the conditional Fourier compiler in
\cite{A2R33}, and this uniform extremal theorem have separate, explicit
hypotheses and conclusions. No implication between them is used without
proof.
